%---------- Inleiding ---------------------------------------------------------

% TODO: Is dit voorstel gebaseerd op een paper van Research Methods die je
% vorig jaar hebt ingediend? Heb je daarbij eventueel samengewerkt met een
% andere student?
% Zo ja, haal dan de tekst hieronder uit commentaar en pas aan.

%\paragraph{Opmerking}

% Dit voorstel is gebaseerd op het onderzoeksvoorstel dat werd geschreven in het
% kader van het vak Research Methods dat ik (vorig/dit) academiejaar heb
% uitgewerkt (met medesturent VOORNAAM NAAM als mede-auteur).
% 

\section{Inleiding}%
\label{sec:inleiding}

Moderne voertuigen genereren enorme hoeveelheden data die het potentieel hebben om bestuurders en eigenaars veel dieper inzicht te geven in de prestaties en gezondheid van hun wagen. Toch blijft deze informatie vaak versnipperd, moeilijk toegankelijk of technisch complex om te interpreteren. In deze bachelorproef richt ik mij op het ontwikkelen van een oplossing die deze technische voertuigdata op een duidelijke en performante manier visualiseert, specifiek voor autoliefhebbers en technisch geïnteresseerde autobezitters die hun auto beter willen begrijpen en nauwkeuriger willen opvolgen.

Het thema wordt gekaderd binnen de bredere evolutie naar data-gedreven voertuigen, waarbij real-time gegevens zoals motortoerental, vermogen, brandstofverbruik of sensordiagnoses steeds vaker beschikbaar zijn via OBD-interfaces, API’s of telematica-systemen. De uitdaging bestaat er echter in om deze data zo te verwerken en te presenteren dat gebruikers er daadwerkelijk inzichten uit kunnen halen, zonder nood aan diepgaande technische of programmeerkennis.

De concrete probleemstelling luidt als volgt: hoe kan voertuigdata in real-time performant, schaalbaar en begrijpelijk worden gevisualiseerd voor autoliefhebbers die hun voertuig nauwkeuriger willen monitoren en analyseren? Vanuit deze probleemstelling vertrekt de centrale onderzoeksvraag:

“Welke technologieën en ontwerpkeuzes maken het mogelijk om real-time voertuigdata performant, schaalbaar en gebruiksvriendelijk te simuleren, streamen en visualiseren binnen een webapplicatie?”

Om deze hoofdvraag te beantwoorden, worden verschillende deelvragen onderzocht:
\begin{itemize}
    \item Welke real-time streamingtechnologie (WebSockets, MQTT of Server-Sent Events) biedt de beste performantie en stabiliteit voor voertuigdata?
    \item Wat is het verschil in realisme en performantie tussen vaste gesimuleerde data en echte OBD-II data logs?
    \item Hoe schaalt de applicatie wanneer het aantal gelijktijdige voertuigen toeneemt?
    \item Welke visualisatietechnieken zijn het meest geschikt voor continue real-time voertuigdata?
    \item Hoe kan OBD-II-data efficiënt worden uitgelezen via vaste, adaptieve en PID-gebaseerde strategieën?
\end{itemize}

De doelstelling van het onderzoek is tweeledig. Enerzijds wil ik een duidelijk en onderbouwd inzicht bieden in welke technologieën (zoals data-streaming, caching, frameworks voor visualisatie) het meest geschikt zijn voor dit type toepassing. Anderzijds werk ik toe naar een concreet proof-of-concept: een webapplicatie die real-time voertuigdata verwerkt en visualiseert op een manier die aansluit bij de noden van autoliefhebbers. Het eindresultaat moet aantonen dat de gekozen technieken effectief leiden tot een vlotte, begrijpbare en betrouwbare gebruikerservaring.

%---------- Stand van zaken ---------------------------------------------------

\section{Literatuurstudie}%
\label{sec:literatuurstudie}

Het domein van real-time voertuigdata, webgebaseerde visualisaties en OBD-II-diagnostiek bevindt zich op het kruispunt tussen automotive informatics, IoT-architecturen en webperformance. De huidige literatuur toont dat zowel de beschikbaarheid van voertuigdata als de manieren waarop deze data performant verwerkt en gevisualiseerd wordt sterk evolueren. Dit onderzoek sluit aan bij bestaande academische inzichten, maar richt zich op een specifieke open uitdaging: hoe kan real-time voertuigdata op een performante, schaalbare én gebruiksvriendelijke manier beschikbaar worden gesteld aan technisch geïnteresseerde autobezitters, zonder afhankelijk te zijn van complexe hardware of bedrijfsomgevingen.

\subsection{OBD-II als standaard voor voertuigdata}

On-Board Diagnostics (OBD-II) vormt sinds 1996 de universele standaard voor voertuigtelemetrie bij consumentenwagens. De review van \textcite{Michailidis2025} toont aan dat OBD-II een brede reeks parameters beschikbaar maakt, zoals motortoerental, voertuigsnelheid, koelvloeistoftemperatuur, lucht-brandstofverhouding en foutcodes (DTC’s). Deze auteurs benadrukken dat OBD-II toegankelijk is, maar tegelijk beperkt in updatefrequentie en afhankelijk van de gebruikte adapter (bijv. ELM327). Deze technische beperkingen zijn relevant voor dit onderzoek, omdat ze bepalen hoe real-time de data werkelijk kan zijn en welke simulatieparameters realistisch zijn.

Daarnaast beschrijven \textcite{Oluwaseyi2020} hoe OBD-II functioneert als digitaal diagnosetool, met een focus op de structuur van PIDs (Parameter IDs) en de uitdagingen rond datakwaliteit en betrouwbaarheid van goedkope adapters. Dit biedt een goede basis voor het deel van dit onderzoek waarin de overgang van gesimuleerde naar echte OBD-II-data wordt onderzocht.

Verder tonen werken zoals \textcite{Dabarera2022} en \textcite{Madushan2023} hoe OBD-II vandaag wordt gecombineerd met IoT-platformen, waarbij data via protocollen zoals MQTT wordt doorgestuurd naar cloudplatformen. Deze studies onderbouwen de haalbaarheid van continue real-time dataflow vanuit voertuigen naar webapplicaties.

\subsection{Real-time data-architecturen en performantie}

Informatie over performante verwerking van real-time data komt vooral uit IoT- en big-dataonderzoek. \textcite{Ferhat2024} beschrijven hoe een voertuigvolgsysteem migratie naar een big-data-omgeving vereist om grote datavolumes efficiënt te verwerken. Hun conclusie is dat traditionele relationele systemen minder geschikt zijn voor real-time data streams, terwijl gedistribueerde systemen met event-driven architecturen betere schaalbaarheid en performantie leveren. Dit is relevant voor de backendarchitectuur van de webapplicatie in deze bachelorproef.

Ook \textcite{Perez-Gonzalez2025} beschrijven een volledig real-time telemetriesysteem voor elektrische racewagens. Hun implementatie toont dat performantie vooral afhankelijk is van:

\begin{itemize}
    \item Sampling rate van sensoren
    \item Datacompressie en transmissiemethode
    \item Efficiënt updaten van dashboards (event-driven i.p.v. polling)
    \item Client-side rendering performance
\end{itemize}

\subsection{Visualisatie van real-time data}

Visualisatie vormt een cruciale component in het interpreteren van telemetrie. \textcite{VandenHautte2020} presenteren een dynamisch dashboarding-framework dat sensordata automatisch detecteert en verwerkt tot real-time visualisaties. Ze benadrukken dat dashboards niet alleen correct moeten zijn, maar vooral responsief, lichtgewicht en geoptimaliseerd voor continue updates.

\subsection{Simulatie van voertuigdata}

Omdat OBD-II-data beperkt is door updatefrequentie en hardwarevariatie, kiezen veel studies voor simulatie als testbasis. In de literatuur rond IoT-telemetrie tonen \textcite{Perez-Gonzalez2025} en \textcite{Madushan2023} dat simulatie cruciaal is om datastromen te valideren nog vóór fysieke sensoren worden aangesloten.

Simulatie wordt vooral ingezet om:
hoge updatefrequenties te testen (bijv. 30–60 updates per seconde),
stress-testing van dashboards te doen,
verschillende voertuigscenario’s (acceleratie, temperatuurverhoging, foutcodes) te reproduceren.

\subsection{Vergelijkende evaluatie als onderzoeksmethode}

Hoewel bestaande studies aantonen dat real-time voertuigdata technisch haalbaar is, blijft er een duidelijk tekort aan vergelijkend onderzoek tussen verschillende streaming-, simulatie- en visualisatie-aanpakken binnen eenzelfde architectuur. Veel implementaties focussen op één technologie of use-case, zonder systematische evaluatie van alternatieven.

Deze bachelorproef richt zich daarom expliciet op een vergelijkend onderzoek, waarbij meerdere technische keuzes binnen dezelfde webapplicatie worden geïmplementeerd en geëvalueerd op basis van performantie, schaalbaarheid en gebruikservaring. Hierdoor wordt niet enkel een werkend systeem opgeleverd, maar ook inzicht verworven in de impact van ontwerpbeslissingen op real-time voertuigdataverwerking.


% Voor literatuurverwijzingen zijn er twee belangrijke commando's:
% \autocite{KEY} => (Auteur, jaartal) Gebruik dit als de naam van de auteur
%   geen onderdeel is van de zin.
% \textcite{KEY} => Auteur (jaartal)  Gebruik dit als de auteursnaam wel een
%   functie heeft in de zin (bv. ``Uit onderzoek door Doll & Hill (1954) bleek
%   ...'')

%---------- Methodologie ------------------------------------------------------
\section{Methodologie}%
\label{sec:methodologie}

Het doel van deze bachelorproef is het ontwerpen, implementeren en evalueren van een performante webapplicatie voor real-time voertuigdata. De methodologie is expliciet vergelijkend van aard en combineert literatuuronderzoek, simulatie, experimentele implementatie en performantie-evaluatie.

In de eerste fase wordt een literatuurstudie uitgevoerd naar OBD-II-standaarden, beschikbare voertuigparameters en technische beperkingen van real-time data-acquisitie. Hierbij wordt onderzocht welke parameters relevant zijn voor real-time visualisatie en welke beperkingen bestaan op vlak van updatefrequentie, betrouwbaarheid en hardware-afhankelijkheid. Deze analyse vormt de basis voor zowel de simulatie als de latere integratie van echte data.

In de tweede fase wordt een simulatieomgeving ontwikkeld waarin voertuigdata op verschillende manieren wordt gegenereerd. Hierbij worden twee simulatie-aanpakken met elkaar vergeleken:
\begin{itemize}
    \item vaste, handmatig gegenereerde datasets met vooraf gedefinieerde rijscenario’s
    \item echte OBD-II-data logs afkomstig van voertuigmetingen
\end{itemize}
Beide aanpakken worden geëvalueerd op realisme, flexibiliteit en performantie binnen de webapplicatie.

Vervolgens wordt de streaminglaag van de applicatie onderzocht. Hiervoor worden meerdere real-time streamingtechnologieën geïmplementeerd en vergeleken, waaronder WebSockets, MQTT en Server-Sent Events. De vergelijking gebeurt op basis van latency, stabiliteit, schaalbaarheid en implementatiecomplexiteit.

In de vierde fase wordt de schaalbaarheid en flexibiliteit van de applicatie geëvalueerd door verschillende voertuigtypes te simuleren, zoals elektrische auto’s en benzinewagens. Hierbij wordt onderzocht welke aanpassingen nodig zijn om de datastromen en parameters van elk type correct te verwerken en visualiseren, en hoe de applicatie stabiel blijft bij hogere updatefrequenties of complexere datastructuren.

Daarnaast wordt een experimentele vergelijking uitgevoerd tussen verschillende visualisatietechnieken voor real-time voertuigdata, zoals line charts, gauges en heatmaps. Deze visualisaties worden geëvalueerd op leesbaarheid, reactietijd en geschiktheid voor continue updates.

Tot slot wordt onderzocht hoe voertuigdata efficiënt kan worden uitgelezen. Hierbij worden drie data-acquisitiestrategieën geïmplementeerd en vergeleken:
\begin{itemize}
    \item vaste polling met een constante updatefrequentie
    \item adaptieve polling gebaseerd op rijgedrag en voertuigstatus
    \item PID-geprioriteerde polling waarbij kritische parameters vaker worden uitgelezen
\end{itemize}

De resultaten van deze experimenten worden geanalyseerd en gebruikt om aanbevelingen te formuleren voor performante en schaalbare real-time voertuigdata-applicaties.


\subsection{Tijdsplanning}

\begin{itemize}
    \item Literatuurstudie en data-analyse (4 weken): overzicht van OBD-II parameters en technische beperkingen.
    \item PoC simulatie (3 weken): simulatie-backend en webdashboard.
    \item Integratie echte OBD-II-data (4 weken): werkende data-acquisitie module en integratie in de applicatie.
    \item Performance-optimalisatie (3 weken): geoptimaliseerde real-time webapplicatie en performance rapport.
    \item Evaluatie en documentatie (2 weken): volledige bachelorproef en technische documentatie.
\end{itemize}

Door deze methodologie wordt zowel een concreet technisch resultaat behaald als diepgaande kennis over OBD-II-data, real-time dataverwerking en webvisualisatie ontwikkeld, waardoor de bachelorproef voldoende IT-diepte en praktische relevantie krijgt.

%---------- Verwachte resultaten ----------------------------------------------
\section{Verwacht resultaat, conclusie}%
\label{sec:verwachte_resultaten}

\subsection{Verwachte Resultaten en Meerwaarde}
Het voornaamste resultaat van deze bachelorproef is een werkende webapplicatie die real-time voertuigdata kan verwerken en visualiseren, zowel vanuit een simulatie als vanaf echte OBD-II apparaten. De applicatie zal dynamische dashboards bevatten met meerdere grafieken die inzicht geven in verschillende voertuigparameters, zoals snelheid, toerental, brandstofverbruik, motortemperatuur en acceleratie.

\subsection{Verwachte technische resultaten}

\begin{itemize}
    \item Een performante backend die data van meerdere voertuigen tegelijk kan verwerken zonder merkbare vertraging
    \item Real-time updates in de frontend via streamingtechnologieën zoals WebSockets of MQTT
    \item Flexibele dashboards die eenvoudig kunnen worden uitgebreid met nieuwe parameters
    \item Een gestandaardiseerde simulatieomgeving waarmee toekomstige uitbreidingen of andere voertuigen makkelijk kunnen worden getest
\end{itemize}


\subsection{Meerwaarde voor de doelgroep}
Voor autoliefhebbers en technisch geïnteresseerde autobezitters biedt deze bachelorproef verschillende concrete voordelen:

\begin{itemize}
    \item \textbf{Inzicht in voertuigprestaties:} Door real-time data visualisatie krijgen gebruikers een dieper begrip van hoe hun voertuig functioneert en waar verbeteringen mogelijk zijn.
    \item \textbf{Diagnose en preventie:} Onregelmatigheden in motor- of rijgedrag kunnen sneller worden opgespoord, wat preventief onderhoud vergemakkelijkt.
    \item \textbf{Gebruiksvriendelijkheid en toegankelijkheid:} In tegenstelling tot standaard dashboards of fabrikant specifieke tools is de applicatie open, flexibel en gericht op de eindgebruiker die technisch geïnteresseerd is.
    \item \textbf{Basis voor verdere ontwikkeling:} De simulatieomgeving en PoC vormen een fundament voor toekomstige uitbreidingen, zoals koppeling met machine learning-algoritmes voor predictive maintenance of rijgedragsanalyse.
\end{itemize}

Kortom, de bachelorproef levert zowel een technisch product op als waardevolle inzichten voor de doelgroep, waardoor zij hun voertuigen beter kunnen begrijpen, bewaken en optimaliseren.

